% ============================================================
% ch00_intro.tex — EXPANDED
% Introduction
% Part of: Tetrad-Based Departure Decomposition (Jiwon Park)
% Session S11 draft — 2026-03-30
% ============================================================

\chapter{Introduction}
\label{ch:intro}


% ============================================================
% §0.1 Observational motivation
% ============================================================

The standard model of cosmology rests on the Friedmann
equation
\begin{equation}\label{eq:Friedmann1}
H^2 = \frac{8\pi G}{3}\,\rho
  - \frac{k}{a^2} + \frac{\Lambda}{3}\,,
\end{equation}
which describes a spatially homogeneous and isotropic
universe whose geometry is determined by the energy density
$\rho$, the spatial curvature $k$, and the cosmological
constant $\Lambda$.  The FLRW framework, built on this
equation, has survived five decades of increasingly precise
tests.

Yet a persistent discrepancy has emerged.  Independent
measurements of the large-scale matter dipole---from CatWISE
mid-infrared quasars, NVSS and RACS radio continuum surveys,
and CosmicFlows-4 peculiar velocities---each report an excess
over the kinematic prediction of the FLRW model.  The
combined significance now exceeds $5\sigma$ across multiple
independent analyses and wavelengths
(Chapter~\ref{ch:dipole}).  The anomaly is not a marginal
tension in a single dataset: it is a multi-probe, multi-wavelength
signal that challenges the foundational assumption of
statistical isotropy.

The question that motivates this work is: \emph{how large is
the departure from FLRW, and what does it constrain?}

To answer this question, we need a framework that quantifies
FLRW departure in a covariant, assumption-explicit manner.
The Maartens--Ellis--Stoeger (MES) kinematic bound
hierarchy~\cite{MES1995a,StoegerME1995} provides the
algebraic foundation: upper bounds on the shear, vorticity,
and four-acceleration relative to the expansion rate, derived
from the observed CMB multipoles.  We extend this hierarchy
into a four-dimensional constraint space through a single
algebraic identity---the \emph{master departure identity}---that
connects all kinematic sources of anisotropy into a scalar
measure of FLRW departure.


% ============================================================
% §0.2 The departure parameter
% ============================================================

The master departure identity is
\begin{equation}\label{eq:intro-MES}
x_C = \Sigstd - \Wstd + \Otilt + \Okaniso\,,
\end{equation}
where $\Sigstd$ is the standardised shear parameter, $\Wstd$
the vorticity parameter, $\Otilt$ the tilt contribution
(from the bulk velocity of the matter frame relative to the
geometry frame), and $\Okaniso$ the anisotropic curvature
contribution.  Each term is a dimensionless ratio,
normalised by the expansion rate $\Theta$, so that
$x_C = 0$ recovers exact FLRW symmetry.

The identity is exact within the Bianchi symmetry class for a
chosen perfect-fluid equation of state.  It is not ``model
independent'' in the sense of holding for arbitrary matter
content, but it is \emph{assumption-explicit}: every
approximation enters at a localised layer in the inference
chain (algebraic identity, bound application, or
observational connection), and the three-layer structure
$x \to Q \to \Pi$ (Chapter~\ref{ch:framework}) makes the
interpretive dependencies auditable.

The departure parameter connects to observables through
calibrated transfer functions.  The shear parameter $\Sigstd$
maps to the CMB quadrupole power $D_2$ through a
species-resolved transfer
(Chapter~\ref{ch:teff}), with
\begin{equation}\label{eq:H-ratio-defect}
\frac{H_\theta}{H_{\rm FLRW}}
= 1 + \frac{1}{2}\,\Sigstd
  + \mathcal{O}\big((\Sigstd)^2\big)
\end{equation}
quantifying the departure of the expansion rate from its
FLRW value.  The deceleration parameter receives a
correction
\begin{equation}\label{eq:delta-q-exact}
\delta q
= \frac{3}{2}\,(1 + w_{\rm eff})\,\Sigstd
  + \mathcal{O}\big((\Sigstd)^2\big)\,,
\end{equation}
and the Jeans length is modified to
\begin{equation}\label{eq:lambda-J-defect}
\lambda_J^{\rm defect}
= \lambda_J^{\rm FLRW}
  \left(1 + \tfrac{1}{2}\,\Sigstd\right)\,.
\end{equation}
For a tilted observer with rapidity $\beta$, the measured
Hubble rate becomes
\begin{equation}\label{eq:tilted-H-tsagas}
H_{\rm tilt}
= H_{\rm FLRW}
  \left(1 + \tanh^2\!\beta\,\Sigstd
    + O(\beta^4)\right)\,,
\end{equation}
and the local expansion is
\begin{equation}\label{eq:vartheta-estimate}
\vartheta
\approx \sqrt{\frac{8\pi G\rho}{3}}
  \cdot \sqrt{1 + \Sigstd}\,.
\end{equation}
These relations quantify how the departure parameter
propagates into physical observables: the expansion rate,
the deceleration, the Jeans scale, and the tilt-induced
Hubble bias.


% ============================================================
% §0.3 Principal results
% ============================================================

Applying the departure framework to current data---the
Planck CMB quadrupole and octupole, the CatWISE/NVSS/RACS
number-count dipoles, and the CosmicFlows-4 bulk flow---we
report a legacy transfer-conditional Bayes-factor summary for a non-zero
tilt-compatible anomaly:
$\ln \mathcal{B}_{\rm tilt} = +25.29 \pm 0.07$
(Route~B, $n_{\rm live} = 1200$, three seeds).  The
corresponding tilt rapidity is
$\beta = 1.360 \times 10^{-3}$, and the Bayesian filling
fraction---the posterior-mean fraction of the MES bound
occupied by the data---is
$\mathcal{F}_{\rm Bayes} = 0.093 \pm 0.025$.

The Phase~1.0 Bianchi Boltzmann solver
(Chapter~\ref{ch:results}) provides the first
species-resolved decomposition of the CMB quadrupole
power $D_2$ within the departure parameter framework,
discovering that neutrinos
contribute $98.4\%$ of $D_2$ through the
Thomson-suppression asymmetry between free-streaming
neutrinos and collision-damped photons.  The massive neutrino
sector is then revisited as an explicit systematic, rather
than a fixed massless idealisation, in
Section~\ref{sec:massive-neutrino-systematic} of
Chapter~\ref{ch:error-hierarchy}.

We emphasise at the outset that the framework quantifies
departure from isotropy under the adopted summary-statistic
likelihood; it does not, by itself, constitute evidence for
anisotropic spatial geometry.  Source identifiability fails
at present: the scalar-amplitude likelihood cannot
distinguish kinematic, geometric, and systematic origins
(geometry false-positive rate $22\%$, minimum
correct-classification rate $68\%$), and we therefore
prohibit geometry claims until direction-dependent data
become available.


% ============================================================
% §0.4 Contributions
% ============================================================

The original contributions of this work are:

\begin{itemize}

\item The master departure identity
  (Eq.~\ref{eq:intro-MES}), extending the MES bound
  hierarchy from one-dimensional scalar bounds to a
  four-dimensional constraint space, derived for all nine
  Bianchi types.

\item A three-layer inference chain $x \to Q \to \Pi$ that
  maps the algebraic departure parameter to a Bayesian
  posterior through the pushforward defect, the occupancy
  ratio, and the exceedance probability.

\item A species-resolved multi-fluid framework
  (Chapter~\ref{ch:framework},
  Section~\ref{sec:multi-fluid}) that assigns each
  cosmological species its own effective-temperature
  manifold, reconstruction map, and tangency diagnostic
  within a single common frame.

\item The multi-species $\Teff$ reduction
  (Chapter~\ref{ch:teff},
  Section~\ref{sec:multi-species-teff}), including the
  FLRW monopole system, the Bianchi quadrupole hierarchy,
  and the neutrino dominance theorem ($N_2/F_2 = 28$,
  $f_\nu = 98.4\%$).

\item A proof that the $\Teff$ closure is immune to the
  Gnedin--Katz moment-closure instability on two
  independent grounds (Section~\ref{sec:gnedin-katz}).

\item The identification of the $\Teff$ tangency diagnostic
  $D_{s,\geq 2}$ with the Gorban--Karlin defect of
  invariance (Section~\ref{sec:invariant-manifold}),
  connecting the cosmological closure to the invariant
  manifold theory of kinetic equations.

\item A Phase~1.0 Bianchi Boltzmann solver with 18
  legacy implementation modules, 331 automated tests, and
  AniCLASS cross-validation at 15 grid points
  (Chapter~\ref{ch:results}).

\item A legacy transfer-conditional Bayes-factor summary
  $\ln \mathcal{B}_{\rm tilt} = +25.29 \pm 0.07$ for
  non-zero tilt, with a seven-channel pseudo-likelihood
  spanning CMB, number counts, and peculiar velocities.

\item A five-level error measurement hierarchy
  (Chapter~\ref{ch:error-hierarchy}) that quantifies the
  solver accuracy from the state space to the observable
  level, confirming that the $\Teff$ closure is not the
  bottleneck.

\item A detailed assessment of the post-Clarkson--Maartens
  mechanism taxonomy: the Tsagas fast-growth mechanism
  is killed; the khronon field and the Krishnan tilt
  instability are the sole surviving tilted-Bianchi
  mechanisms (Chapter~\ref{ch:discussion},
  Section~\ref{sec:disc-mechanisms}).

\end{itemize}


% ============================================================
% §0.5 Scope and limitations
% ============================================================

The framework is subject to three structural limitations
that we state explicitly.

First, the Bianchi symmetry class assumes spatial
homogeneity.  The departure parameter $x_C$ constrains the
global anisotropy but cannot address spatially
inhomogeneous backreaction.  The connection to the Buchert
spatial-averaging framework is discussed in
Chapter~\ref{ch:future}
(Section~\ref{sec:backreaction-programme}), where we
identify the structural mismatch between the tensorial
Bianchi curvature and the scalar Buchert average as the
principal obstacle to a complete dictionary.

Second, the inference pipeline uses a summary-statistic
likelihood (seven channels compressed to scalar amplitudes),
not a pixel-level analysis of the CMB sky.  The compression
discards directional information that could break the
equivalence classes among Bianchi types
(Chapter~\ref{ch:robustness}).  The direction-dependent
programme (Section~\ref{sec:phase30}) is designed to
recover this information.

Third, the $\Teff$ closure provides a monitored
finite-dimensional reduction of the infinite-dimensional
Boltzmann hierarchy, not a closed dynamical system with
proven global-in-time stability.  The tangency diagnostic
$D_{s,\geq 2}$ certifies instantaneous adequacy;
long-time accuracy requires separate verification
(Chapter~\ref{ch:error-hierarchy}).


% ============================================================
% §0.6 The Ehlers--Geren--Sachs theorem
% ============================================================

\section{The Ehlers--Geren--Sachs theorem and isotropy}
\label{sec:intro-EGS}

The Ehlers--Geren--Sachs (EGS) theorem establishes that if
all fundamental observers in an expanding universe measure
an exactly isotropic CMB, the spacetime is exactly FLRW.
The almost-EGS extensions of Stoeger, Maartens \&
Ellis~\cite{StoegerME1995} allow for small anisotropies
consistent with observations: if the CMB is nearly isotropic
(to order $\varepsilon$), the spacetime is nearly FLRW (to
order $\varepsilon$, with possible cumulative secular
growth).

The dipole anomaly does not violate the EGS theorem
directly---the CMB is isotropic to $10^{-5}$ in the CMB rest
frame---but it implies that the matter rest frame is
displaced from the CMB rest frame by more than the
$\Lambda$CDM kinematic prediction.  The departure framework
quantifies this displacement through the tilt parameter
$\Otilt$, which measures the energy flux carried by the
relative motion of matter and radiation.  The EGS theorem
guarantees that if this displacement vanishes, the FLRW
model is recovered; the departure framework tracks what
happens when it does not.


% ============================================================
% §0.7 Thesis outline
% ============================================================

\section{Outline of the thesis}
\label{sec:outline}

The remainder of this thesis is organised as follows.

Chapter~\ref{ch:dipole} reviews the observational evidence
for the cosmic dipole anomaly: the number-count dipole
excess, the peculiar velocity surveys, the statistical
methodology, and the near-term experimental outlook.

Chapter~\ref{ch:framework} develops the mathematical
framework: the 1+3 covariant decomposition, the standardised
kinematic variables, the master departure identity, the
three-layer inference chain, the three-frame problem, and the
multi-fluid species-resolved extension.

Chapter~\ref{ch:bianchi-bounds} derives the MES bounds for
all nine Bianchi types and constructs the type-by-type
constraint atlas.

Chapter~\ref{ch:teff} presents the effective temperature
formalism: the $\Teff$ ansatz, the $\Theta^4$ moment map,
the nonlinear corrections, the departure propagation, the
multi-species reduction with the neutrino dominance theorem,
the Gnedin--Katz response, and the nonperturbative regime.

Chapter~\ref{ch:pipeline} describes the observational
inference pipeline: the three-code architecture (BASS, HTT,
MIO), the data sources, the code modules, the Bianchi
Boltzmann solver architecture, and the validation protocol.

Chapter~\ref{ch:results} presents the numerical results:
the Bayesian evidence, the filling fraction, the
equivalence-class analysis, the channel decomposition, and
the species-resolved solver validation.

Chapter~\ref{ch:robustness} analyses robustness through
null recovery, channel ablation, prior sensitivity,
posterior predictive checks, and model identifiability audits.

Chapter~\ref{ch:discussion} discusses the implications:
the contrastive analysis of principal results, the Dipole
Cosmological Principle, the Tsagas programme, mechanism
attribution after Clarkson--Maartens, and multi-fluid
implications.

Chapter~\ref{ch:future} outlines future directions: testable
predictions, the experiment-by-experiment assessment, the
anisotropic Boltzmann solver programme, and the backreaction
roadmap.

Chapter~\ref{ch:error-hierarchy} defines the five-level error
measurement hierarchy for the solver programme.
